\section*{Bibliography}
\phantomsection
\addcontentsline{toc}{section}{Bibliography}

\begin{enumerate}
\def\labelenumi{(\Roman{enumi})}
\item
  O. NEUGEBAUER, Lectures on the History of Ancient Mathematical Sciences, Vol. I: Pre-Greek Mathematics, Berlin (Springer), 1934.
\item
  J. TROPFKE, History of Elementary Mathematics, vols. IV–VI, Berlin–Leipzig (de Gruyter), 1923–24.
\item
  Euclid's Elements, 5 vols., ed. J. L. Heiberg, Leipzig (Teubner), 1883–88.
\end{enumerate}

(III bis) T. L. HEATH, The thirteen books of Euclid's Elements\ldots, 3 vols., Cambridge, 1908.

\begin{enumerate}
\def\labelenumi{(\Roman{enumi})}
\setcounter{enumi}{3}
\item
  T. L. HEATH, Apollonius of Perga, Treatise on conic sections, Cambridge (Univ. Press), 1896.
\item
  Ptolemaei Cl. Opera, ed. J. L. Heiberg, 2 vols., Leipzig (Teubner), 1898–1903.
\item
  G. DESARGUES, WORKS\ldots vol. I, Paris (Leiber), 1864: Brouillon project d'une atteinte aux événements des rencontres d'un cône avec un plan, pp. 103-230.
\item
  P. FERMAT, Œuvres, vol. I, Paris (Gauthier-Villars), 1891: a) Ad locos planos et solidos Isagoge, pp. 91-110 (French translation, ibid., vol. III, pp. 84-101); b) Isagoge ad locos ad superficiem, pp. 111-117 (French translation, ibid., vol. III, pp. 102-108).
\item
  L. EULER: a) Introductio in Analysin Infinitorum (Opera Omnia (1), vol. IX, Zürich-Leipzig-Berlin (O. Füssli and B. G. Teubner), 1945); b) Theoria motus corporum solidorum seu rigidorum (Opera Omnia (2), vol. III, Zürich-Leipzig-Berlin (O. Füssli and B. G. Teubner), 1948); c) Problema algebraicum ob affections prorsus singulares memorabile (Opera Omnia, (1), vol. VI, Leipzig-Berlin (Teubner), 1921, pp. 287-315); d) Formulae generales pro translatione quacunque corporum rigidorum, Novi Comm. Acad. Sc. imp. Petrop., vol. XX (1776), pp. 189-207; e) De centro similitudinis. (Opera Omnia (1), vol. XXVI, Zürich (O. Füssli), 1956, pp. 276-285).
\item
  J. L. LAGRANGE, Œuvres, Paris (Gauthier-Villars), 1867-1892: a) Researches on the method of maxima and minima, vol. I, pp. 3–20; b) Researches in arithmetic, vol. III, pp. 695–795.
\item
  G. MONGE, Descriptive Geometry, Paris, 1798.
\item
  C. F. GAUSS, Works: a) Disquisitiones arithmeticae, vol. I, Göttingen, 1870; b) Mutations of Space, vol. VIII, Göttingen, 1900, pp. 357–362.
\item
  J.-V. PONCELET, Treatise on the Projective Properties of Figures, vol. I, 2nd ed., Paris (Gauthier-Villars), 1865.
\item
  A. F. Möbius, Collected Works, Leipzig (Hirzel), 1885–87: a) The Barycentric Calculus, vol. I, pp. 1–388; b) On a Special Kind of Dual Relations between Figures in Space, vol. I, pp. 489–515 (= Crelle's Journal, vol. X, 1833); c) On a New Treatment of Analytic Spherics, vol. II, pp. 1–54.
\item
  A.-L. CAUCHY: a) Lessons on the applications of infinitesimal calculus to geometry (Complete Works, (2), vol. V, Paris (Gauthier-Villars), 1903); b) On the equation by means of which one determines the secular inequalities of the planets (Complete Works, (2), vol. IX, Paris (Gauthier-Villars), 1891, pp. 174–195).
\item
  M. CHASLES: a) Note on the general properties of the system of two bodies, Bull. de Férussac, vol. XIV (1830), pp. 321–326; b) Historical survey of the origin and development of methods in geometry, Brussels, 1837.
\item
  E. GALOIS, Mathematical Works, Paris (Gauthier-Villars), 1897.
\item
  W. R. HAMILTON, Lectures on Quaternions, Dublin, 1853.
\item
  A. CAYLEY, Collected Mathematical Papers, Cambridge, 1889-1898: a) On certain results relating to quaternions, vol. I, p.~123-126 (= Phil. Mag., 1845); b) Sur les déterminants gauches, vol. I, p.~410-413 (= J. de Crelle, vol. XXXVIII (1848)); c) Recherches ultérieures sur les déterminants gauches, vol. II, p.~202-215 (= J. de Crelle, vol. L (1855)); d) A sixth memoir on quantics, vol. II, p.~561-592 (= Phil. Trans., 1859).
\item
  C. G. J. JACOBI, Gesammelte Werke, Berlin (G. Reimer), 1881-1891: a) On the Pfaffian method\ldots vol. IV, pp. 17–29; b) On a fundamental algebraic theorem and its applications, vol. III, pp. 593–598.
\item
  J. J. SYLVESTER, Collected Mathematical Papers, vol. I, Cambridge, 1904: A demonstration of the theorem that every homogeneous quadratic polynomial is reducible by real orthogonal substitution to the form of a sum of positive and negative squares, pp. 378–381 (= Phil. Mag., 1852).
\item
  E. LAGUERRE, Works, vol. II, Paris (Gauthier-Villars), 1905.
\item
  C. HERMITE, Works, vol. I, Paris (Gauthier-Villars), 1905: On the theory of quadratic forms, pp. 200–263 (= J. de Crelle, vol. XLVII (1854)).
\item
  K. G. V. von STÄUDT, Contributions to the Geometry of Position, Nuremberg, 1856.
\item
  E. BELTRAMI: a) Essay on the interpretation of non-Euclidean geometry, Giorn. di Mat., vol. VI (1868), pp. 284–312; b) Fundamental theory of spaces of constant curvature, Ann. di Mat. (2), vol. II (1868–69), pp. 232–255.
\item
  F. KLEIN, Collected Mathematical Works, vol. I, Berlin (Springer), 1921: a) On the so-called non-Euclidean geometry, p.~254-305 (= Math. Ann., vol. IV (1871)); b) Comparative considerations on recent geometrical research, p.~460-497 (= Math. Ann., vol. XLIII (1893)).
\item
  C. JORDAN, Treatise on substitutions and algebraic equations, Paris (Gauthier-Villars), 1870.
\item
  G. FROBENIUS; a) On linear substitutions and bilinear forms, J. de Crelle, vol. LXXXIV (1878), p.~1-63; b) Theory of linear forms with integer coefficients, J. de Crelle, vol. LXXXVI (1879), p.~146-208.
\item
  W. K. CLIFFORD, Mathematical Papers, London (Macmillan), 1882: a) On the classification of geometric algebras, p.~397-401; b) Applications of Grassmann's extensive algebras, p.~266-276 (= Amer. Journ. of Math., vol. I (1878)).
\item
  R. LIPSCHITZ, Investigations on the sums of squares, Bonn, 1886.
\item
  D. HILBERT, Foundations of Geometry, Leipzig (Teubner), 1899.
\item
  E. WITT, Theory of quadratic forms in arbitrary fields, J. de Crelle, vol. CLXXVI (1937), p.~31-44.
\item
  E. CARTAN, Lectures on the theory of spinors, Actual. Sci. et Industr., nos. 643 and 701, Paris (Hermann), 1938.
\item
  C. L. SIEGEL, Symplectic Geometry, Amer. Journ. of Math., vol. LXV (1943), p.~1-86.
\item
  C. CHEVALLEY, The algebraic theory of spinors, New York (Columbia Univ. Press), 1954.
\item
  M. EICHLER, Quadratic forms and orthogonal groups, Berlin-Göttingen-Heidelberg (Springer), 1952.
\end{enumerate}

